\section{Remove the decoder's second source language}

\paragraph{Recommendation.}
Index the decoder roles directly by the existing register program, rather than
by the separately numbered \texttt{DecoderCertificate.Node} array. This is a
representation deletion, not another graph minimization or a new source proof.
The current correspondence is already kernel-proved; no trust gap is alleged.
This simplifies the separately checked prefix component and its future
integration, not the current end-to-end endpoint closure: production still
uses the fixed-fuel checker and all 6,176 counted guarded outcomes.

\paragraph{Exact pinned finding.}
At commit \texttt{cbf94cb552a3b49509e0eee80e7455d4eb5974c6},
\texttt{K215Decoder.sourceNodes} has 202 entries: 179 forks and 23 primitive
leaves. Its forks, enumerated in their existing order, are precisely the 179
nodes of \texttt{KSource215.p215}. The decoder root 201 maps to program root 178.
For the map $\rho$ from alternate nodes to register edges,
\[
 \rho(v)=
 \begin{cases}
   \operatorname{node}(j),&v\text{ is the }j\text{-th fork},\\
   \operatorname{primitive}(o,a),&v\text{ is the leaf }(o,a),
 \end{cases}
\]
the audit checks every local square
\[
 \mathrm{sourceNodes}[v]=\operatorname{fork}(l,r)
 \quad\Longrightarrow\quad
 \mathrm{p215.nodes}(j)=(\rho(l),\rho(r)).
\]
All 179 squares pass. All 468 role branches retarget: 371 internal and 97
terminal. The 234 roles and 281 physical blocks are preserved, not identified
with source nodes or merged with each other.

\subsection{The layer that can disappear}
The current root theorem starts with \texttt{RegisterPath.walk}, converts through
\texttt{twalk} and \texttt{LeafAt}, rewrites by
\texttt{source\_tree\_eq}, and reconstructs a
\texttt{SourcePath sourceNodes} witness. That witness is then interpreted through
the role and block certificate. The separate source tree is unfolded with fuel
203, although the actual program already supplies its own decoder and fuel.

Instead, a role's source field should be an actual \texttt{Program} node.
An internal link still stores its target-role and block indices. A terminal
link stores its physical action, but no separately numbered primitive leaf.
The selected logical operation is read directly from the actual program's
left or right edge. The new checker demands exactly:
\begin{itemize}
\item Internal source edge: the actual edge is \texttt{node child}; the target
  role names \texttt{child}; its short suffix, physical state and referenced
  block satisfy the current window and endpoint equalities. Preserve
  decreasing labels on internal source references.
\item Terminal source edge: the actual edge is not a node, and the stored
  physical action equals the real table row at the role's state and read bit.
  The terminal operation and argument remain those of the actual source edge.
\item Root and domain: the first role names \texttt{p.root}, state 84 and the
  initial marker suffix; role source labels are below \texttt{p.nnodes}.
\end{itemize}

The full consumed prefix remains part of the routed-entry predicate. The
physical state or its two-bit history is not a substitute for a role path.
A terminal action/table-row equality is not an opcode-semantics proof;
residual primitive and return contracts remain mandatory.

\subsection{One direct walk induction}
The required semantic replacement is a single induction on a successful
\texttt{walk} equation. A convenient invariant uses an absolute tape offset
$a$, consumed position $i$, and final depth $d$. Assume
\[
 \operatorname{walk}(p,f,n,P,i)=\operatorname{some}(e,d),
 \qquad t(a+j)=\operatorname{bit}(P,j)\quad(i\leq j<d),
\]
and that the checked current role names source node $n$ and its suffix is
present immediately before $a+i$. The theorem constructs a certified block
chain from the role's physical state at head $a+i$ to the terminal primitive
port at head $a+d-1$, preserving the entire arbitrary tape. It also records
the exact source operation $e$ and the complete routed-entry witness.

For a terminal source edge, success gives $d=i+1$ and the internal chain is
empty. For an internal edge, the structural certificate selects the next role
and block. Existing \texttt{before\_snoc}, \texttt{before\_drop},
\texttt{before\_window} and arbitrary-tape block soundness establish the next
boundary. Apply the induction hypothesis at $i+1$ and prepend that block.
There are exactly $d-i-1$ internal links; the unchanged duration facts give
\[
 d-i-1\ \leq\ \mathrm{clock}\ \leq\ 7(d-i-1).
\]
At the root, the interval-bit assumption follows from the existing
\texttt{Reads} hypothesis on \texttt{consume P d}. Using $P$ itself would be
incorrect: the decoder may consume implicit trailing zeros or ignore a
supplied suffix. The theorem is proposition-valued and requires no forbidden
extraction of a computational clock from a general proof.

This proof is proposed, not compiled here. The owned finite audit does not
establish its quantified tape semantics. The existing zipper reflection and
physical block proofs should be reused unchanged.
The minimal prefix-library edit set is \texttt{DecoderCertificate},
\texttt{K215Decoder}, \texttt{DecoderPaths}, \texttt{DecoderSourceBinding},
and \texttt{KDecoderPrefix215}. For this step, \texttt{KSim}, the nine
\texttt{KSites215Batch} modules and \texttt{KSites215} remain unchanged.

\subsection{Freshness, compatibility and limits}
This differs from the already implemented computed-tree change: defining
\texttt{t215 := ptree p215} removed the expanded tree literal but retained
\texttt{sourceNodes}, the second source language and their correspondence.
It also differs from earlier program-first plans that choose a route and clock
before introducing the tape, or use decreasing labels for coverage. Those
ideas can be composed with this representation deletion, but are not repeated
as new results.

The direct root theorem consumes the same \texttt{walk} result used by
\texttt{KSim.Covered}. It does not itself prove total source coverage,
primitive/return contracts, or acceptance by the fixed-fuel Boolean site
checker. Existing production discharges its old checks; a replacement must
preserve all 6,176 counted guarded outcomes. Semantic
reachability must not be used as evidence that fixed-fuel checking succeeds.

The internal routed-entry type changes. Preserve final machine/refutation and
consistency endpoints, and audit internal callers. If exact legacy decoder
types must survive, retain a separate compatibility facade and count it in
the maintenance cost. A generic inverse \texttt{LeafAt}-to-\texttt{twalk}
lemma can retain a leaf-facing wrapper without the alternate source array;
preserving the old exact result type is a different requirement.

\paragraph{Delivered evidence.}
The owned parser verifies exact pinned Git blob identities before reading
the two literal graphs. Its output includes the complete 179-label map,
23 primitive projections and all 234 role relabelings. It checks all 4,751
expanded leaf walks, counts 6,176 decrement-sensitive outcomes, and exercises
unread suffixes and trailing-zero erasure. Swapped fork labels, a changed
primitive and a changed root child are rejected by negative controls. These
are finite source-data checks, not runtime guard checks or a new proof-assistant
receipt. An independent strict-parser review reproduced the counts.

\paragraph{Deliverable boundary.}
\texttt{DirectProgramCertificateDraft.lean} supplies the new finite checker,
literal retargeted roles and proposed kernel-check declarations. It has not
been compiled and supplies no generic direct-walk soundness theorem. Its old
decoder import is a migration convenience for the 281 blocks; production must
first split those blocks into a source-independent module to obtain the
intended dependency deletion.

The old 202-node literal is 3,132 UTF-8 bytes; the current
\texttt{DecoderSourceBinding.lean} is 4,456 bytes. These are gross deletion
candidates, not a net-saving claim. Count the replacement induction, checker,
updated clients and any facade. No proof-object, runtime, memory or build-time
improvement is claimed. No corresponding Rocq representation deletion is
asserted without inspecting an analogous source boundary.
